%%%%%%%%%%%%%%%%%%%%%%%%%%%%%%%%%%%%%%%%%%%%%%%%%%%%%%%%%%%%%%%%%%%%
%% I, the copyright holder of this work, release this work into the
%% public domain. This applies worldwide. In some countries this may
%% not be legally possible; if so: I grant anyone the right to use
%% this work for any purpose, without any conditions, unless such
%% conditions are required by law.
%%%%%%%%%%%%%%%%%%%%%%%%%%%%%%%%%%%%%%%%%%%%%%%%%%%%%%%%%%%%%%%%%%%%

\documentclass[
  digital,     %% The `digital` option enables the default options for the
               %% digital version of a document. Replace with `printed`
               %% to enable the default options for the printed version
               %% of a document.
%%  color,       %% Uncomment these lines (by removing the %% at the
%%               %% beginning) to use color in the printed version of your
%%               %% document
  oneside,     %% The `oneside` option enables one-sided typesetting,
               %% which is preferred if you are only going to submit a
               %% digital version of your thesis. Replace with `twoside`
               %% for double-sided typesetting if you are planning to
               %% also print your thesis. For double-sided typesetting,
               %% use at least 120 g/m² paper to prevent show-through.
  nosansbold,  %% The `nosansbold` option prevents the use of the
               %% sans-serif type face for bold text. Replace with
               %% `sansbold` to use sans-serif type face for bold text.
  nocolorbold, %% The `nocolorbold` option disables the usage of the
               %% blue color for bold text, instead using black. Replace
               %% with `colorbold` to use blue for bold text.
  lof,         %% The `lof` option prints the List of Figures. Replace
               %% with `nolof` to hide the List of Figures.
  lot,         %% The `lot` option prints the List of Tables. Replace
               %% with `nolot` to hide the List of Tables.
]{fithesis4}
%% The following section sets up the locales used in the thesis.
\usepackage{polyglossia}  %% By using `czech` or `slovak` as the
\setmainlanguage{czech} %% main locale instead of `english`, you
%% can typeset the thesis in either Czech or Slovak, respectively.
\setotherlanguages{english} %% The
%% additional keys allow foreign texts to be typeset as follows:
%%
%%   \begin{otherlanguage}{german}  ... \end{otherlanguage}
%%   \begin{otherlanguage}{russian} ... \end{otherlanguage}
%%   \begin{otherlanguage}{czech}   ... \end{otherlanguage}
%%   \begin{otherlanguage}{slovak}  ... \end{otherlanguage}
%%
%% For non-Latin scripts, it may be necessary to load additional
%% fonts:
\newfontfamily\russianfont[Script=Cyrillic,Ligatures=TeX]{PT Serif}
%%
%% The following section sets up the metadata of the thesis.
\thesissetup{
    date        = \the\year/\the\month/\the\day,
    university  = mu,
    faculty     = fi,
    type        = bc,
    department  = Katedra strojového učení a zpracování dat,
    author      = Jiří Kyjovský,
    gender      = m,
    advisor     = {Ing. Matej Lexa, Ph.D.},
    title       = {TODO},
    TeXtitle    = {TODO},
    keywords    = {TODO},
    TeXkeywords = {TODO},
    abstract    = {%
      TODO
    },
    thanks      = {%
      TODO
    },
    bib         = bibliography.bib,
    %% Remove the following line to use the JVS 2018 faculty logo.
    facultyLogo = fithesis-fi,
}
\usepackage{makeidx}      %% The `makeidx` package contains
\makeindex                %% helper commands for index typesetting.
\usepackage[acronym]{glossaries}          %% The `glossaries` package
\renewcommand*\glspostdescription{\hfill} %% contains helper commands
\loadglsentries{example-terms-abbrs.tex}  %% for typesetting glossaries
\makenoidxglossaries                      %% and lists of abbreviations.
%% These additional packages are used within the document:
\usepackage{paralist} %% Compact list environments
\usepackage{amsmath}  %% Mathematics
\usepackage{amsthm}
\usepackage{amsfonts}
\usepackage{url}      %% Hyperlinks
\usepackage{markdown} %% Lightweight markup
\usepackage{listings} %% Source code highlighting
\lstset{
  basicstyle      = \ttfamily,
  identifierstyle = \color{black},
  keywordstyle    = \color{blue},
  keywordstyle    = {[2]\color{cyan}},
  keywordstyle    = {[3]\color{olive}},
  stringstyle     = \color{teal},
  commentstyle    = \itshape\color{magenta},
  breaklines      = true,
}
\usepackage{floatrow} %% Putting captions above tables
\floatsetup[table]{capposition=top}
\usepackage[babel]{csquotes} %% Context-sensitive quotation marks
\begin{document}
%% Uncomment the following lines (by removing the %% at the beginning)
%% and to print out List of Abbreviations and/or Glossary in your
%% document. Titles for these tables can be changed by replacing the
%% titles `Abbreviations` and `Glossary`, respectively.
%% \clearpage
%% \printnoidxglossary[title={Abbreviations}, type=\acronymtype]
%% \printnoidxglossary[title={Glossary}]

\chapter*{Úvod}
\markright{\textsc{Úvod}}
\addcontentsline{toc}{chapter}{Úvod}

TODO na konci bakalarky


% ==================================================================================== %
% ==================================================================================== %


\chapter{Aktuální stav}

Studenti bioinformatiky na FI potřebují ke svému studiu využívat různý
bioinformatický software, stejně tak vyučující pro výuku. Některý software
je snadno přístupný a může být nainstalován přímo z repozitáře
distribuce, která je používána studentem, nebo je dostupný z nějakého jiného
repozitáře jako je například PyPI. Ne všechen potřebný software je takto dostupný
a studenti či vyučující by si museli daný software zkompilovat. Fakulta proto 
provozuje vitruální stroj s operačním systémem Bio-Linux \cite{biolinux_fi},
který byl využíván pro účely výuky, na nějž se mohou studenti připojit a
pracovat tam.

% ------------------------------------------------------------------------------------ %

\section{Seznámení se s problémem}

Bio-Linux, distribuce Linuxu na bázi Ubuntu 14.04 LTS, byl projektem snažícím se poskytnout
potřebnou množinu programů pro uživatele zajímající se o bioinformatiku
\cite{biolinux}. Tato distribuce se bohužel již od roku 2019 nachází ve stavu 
EOL\footnote{EOL znamená End of Life} \cite{biolinux_distrowatch}, čímž vznikla 
nutnost jej nahradit, jelikož provozování softwaru, který je ve stavu EOL, může
vést k nepříjemnostem \cite{eol}. Je tedy potřeba studentům a vyučujícím na FI
poskytnout nějakou základní množinu bioinformatického softwaru, který by byl již
zabalen a připraven k instalaci na zařízení pro potřeby výuky a nahradit tak 
zastaralý Bio-Linux.

Úkolem bakalářské práce tedy bylo navrhnout vhodné řešení k
zabalení softwaru, vytvořit nebo využít nějaké úložiště pro tyto balíčky,
poskytnout uživatelsky přívětivou cestu, jak si balíčky stáhnout a nainstalovat je
na zařízení. Dále bylo potřeba navrhnout celý systém tak, aby do něj mohl být dle
potřeby okamžitě přidán či aktualizován software potřebný pro výuku. Celý systém
by neměl být náročný na údržbu nebo aktualizace a přidávání softwaru by mělo být
jednoduché, aby s tím neměl vyučující potíže. Takto připravený software se pak
navíc může nainstalovat na nový virtuální stroj, který bude fakulta provozovat
jako v případě Bio-Linuxu.

% ------------------------------------------------------------------------------------ %

\section{Ubuntu}

TODO aktualni stav bioinforrmatickych baliku v ubuntu


% ------------------------------------------------------------------------------------ %

\section{Fedora}

TODO aktualni stav bioinforrmatickych baliku ve fedora


\subsection{Copr}

% ------------------------------------------------------------------------------------ %

\section{Conda}

TODO aktualni stav bioinforrmatickych baliku pro condu


\subsection{Conda-forge}


% ==================================================================================== %
% ==================================================================================== %


\chapter{Obecný návrh}

Hlavní motivací při návrhu celého systému pro mě bylo použít co nejvíce fungujících
nástrojů a služeb, spojit je dohromady tak, aby vytvořily funkční celek a soustředit
se co nejvíce na automatizaci a zpříjemnění interakce se systémem pro jeho uživatele.
Celý systém jsem tedy rozdělil do čtyř hlavních komponent:

\paragraph{Build systém}
Jádro celého systému. Služba, která bude schopna sestavit balíček z daného softwaru
a tento balíček vrátit.

\paragraph{Úložiště}
Databáze, která bude obsahovat již sestavené balíčky, vrácené z build systému, které
jsou připraveny k instalaci na zařízení uživatele.

\paragraph{Služba pro interakci správce se systémem}
Pomocí této služby bude správce komunikovat s build systémem a databází v intuitivní
a abstraktní rovině. Detaily build systému nebo databáze zůstanou správci skryty.

\paragraph{Sada nástrojů pro interakci}
Doplňující, již existující nebo mnou vytvořená, množina nástrojů, která bude spojovat
tři výše zmíněné komponenty do funkčního celku.

\bigbreak
Zavedl jsem pohledy správce a uživatele na systém a podle nich jsem postupoval nezávisle
na zvolení build systému nebo úložiště tak, aby bylo těmto pohledům vyhověno a zároveň
abych mohl později zvolit vhodný build systém s úložištěm.

% ------------------------------------------------------------------------------------ %

\section{Pohled uživatele}

Instalace softwaru z pohledu uživatele musí být co nejméně náročná a zároveň
by se neměly na uživatele přenášet žádné složitější úkony, aby byla práce
se systémem co nejjednodušší. V ideálním případě by měl být celý systém
navržen tak, aby si uživatel nainstaloval software pomocí jediného příkazu.

V případě interakce uživatele s virtuálním strojem musí být tento stroj rovnou
připravený k používání poté, co se k němu uživatel připojí, což znamená, že
množina softwaru musí být na virtuálním stroji již předinstalována.

% ------------------------------------------------------------------------------------ %

\section{Pohled správce}

Správce bude mít za úkol přidávat nový software do systému, případně jej
může aktualizovat. Po přidání nového softwaru nebo jeho aktualizace se
spustí série automatizovaných úkonů, které zaručí, že software bude
správně zabalen a bude dostupný uživateli k instalaci.

Aktualizace softwaru by neměla být primární zodpovědnost správce, systém
totiž bude periodicky kontrolovat nové verze softwaru a pokusí se je
sám aktualizovat. V případě neúspěchu bude správce o tomto skutku
kontaktován prostřednictvím emailu.

% ------------------------------------------------------------------------------------ %

\section{Repozitář metadat}

Pro interakci mezi správcem a systémem jsem se rozhodl použít podobný
koncept, který používá Flathub, což je schraňování potřebných metadat
aplikace, jako je například seznam závislostí nebo zdroj upstreamu, v
git repozitáři. Odsud se aplikace automatizovanou cestou díky
poskytnutým metadatům sestaví a odešle na úložiště Flathubu.
Přidání nové aplikace se řeší přes pull request na GitHubu 
a po schválení administrátory se vytvoří nový repozitář v organizaci
\textit{github.com/flathub}
\cite{flatpak_app_creation, flatpak_app_repo, flatpak_app_build}, do
kterého mají administrátoři přístup a mohou daný software dále
aktualizovat přes pull requesty v jejich přiděleném repozitáři pro
danou aplikaci \cite{flatpak_app_update}.

Podobný koncept však není nic nového. Například Fedora používá pro
ukládání metadat také řešení na konceptu gitu a říká jí 
dist-git \cite{dist-git}. Všechna metadata balíčků má uložené ve
své vlastní git službě\footnote{Dostupné na 
\url{https://src.fedoraproject.org/}}. Conda-forge také používá
podobné řešení na GitHubu\footnote{Všechny balíčky jsou dostupné
na \url{https://github.com/orgs/conda-forge/repositories}} jako
Flathub \cite{conda-forge_repo_creation}.


\subsection{Struktura úložiště metadat}
\label{item:2.3.1}

Stejně jako v případě řešení Flathubu jsem se rozhodl ukládat metadata
balíčků na GitHubu, kde správce bude moct pomocí pull requestu přidávat
nové balíčky nebo ty stávající aktualizovat.

Pro jednoduchost jsem nepoužil myšlenku jeden repozitář - jeden balíček,
nýbrž všechny balíčky budou skladovány v jednom velkém repozitáři, kde
každý příslušný balíček bude mít svůj adresář s metadaty. Rozdělení 
každého balíčku do separátního repozitáře je vhodné pro situaci, kde
máme mnoho správců starající se o své balíčky a jiné je nezajímají.
Toto rozdělení zajišťuje oddělené práva přístupu. Naproti tomu v
našem případě se o balíčky bude starat pouze jedna entita a to je
fakulta, takže pro správce bude jednodušší, když bude mít všechny
své balíčky na jednom místě.

\begin{figure}[H]
  \begin{center}
    \includegraphics[width=0.6\textwidth]{images/png/repository_structure.png}
  \end{center}
  \caption{Ukázka struktury GitHub repozitáře}
  \label{fig:repository_structure}
\end{figure}


\subsection{Metadata balíčku}

Jak můžeme vidět na obrázku \ref{fig:repository_structure}, každý
balíček obsahuje soubor \textit{metadata.yaml}, který specifikuje
určitá metadata:

\paragraph{maintainers}
Seznam správců, kteří zodpovídají za konkrétní balíček. Tento seznam
bude u každého správce obsahovat jeho GitHub nickname a email.

\paragraph{url}
URL odkud bude upstream daného balíčku dostupný ke stažení.

\paragraph{version}
Aktuální verze balíčku v době aktualizace nebo přidání. Podle téhle
informace bude systém určovat, jestli je daný balíček aktuální.

\paragraph{sha512}
Checksum, konkrétně typu SHA512, podle kterého si bude systém ověřovat
pravost stáhnutého zdroje.

\paragraph{autoupdate}
Výchozí hodnota nastavena na \textbf{True}. Určuje, jestli se má systém
pokoušet balíček automaticky aktualizovat.


\subsection{Správa balíčků}

Jak již bylo zmíněno v sekci \ref{item:2.3.1}, správce bude interagovat
se systémem pomocí pull requestu. V případě přidání nového balíčku
správce vytvoří v repozitáři nový adresář. Jeho jméno bude použito jako
jméno balíčku. V tomto adresáři správce vyplní potřebná metadata a systém
skrze GitHub Action\footnote{CI$\backslash$CD systém od 
GitHubu\cite{gh_actions}} automaticky otestuje, jestli se balíček podaří
sestavit. Pokud ano, může správce pull request spojit s hlavní větví a 
přes GitHub Actions se balíček pošle do databáze tak, aby byl dostupný 
uživateli k instalaci. V případě aktualizace bude celá situace obdobná,
pouze správce upraví již existující metadata balíčku tak, aby se 
aktualizace správně sestavila.


\subsection{Automatická aktualizace balíčků}

Systém bude periodicky hlídat upstream verzi balíčku (co měsíc/týden/den)
a při zjištění nové verze se provede pokus o aktualizaci. Bohužel je nevýhodné
využít již navrženého schématu k aktualizaci skrze pull request kvůli
přísným limitům GitHubu na vytváření obsahu \cite{gh_secondary_rate_limits}
a velmi omezujícímu způsobu, jak se těmto chybám s limity vyhnout
\cite{gh_avoid_secondary_rate_limits}.

Pokus o aktualizaci se provede vytvořením nové větve a změnou atributu
\textbf{version} v metadatech balíčku na nejnovější verzi. Po aplikování
změn do této větve se nad ní spustí stejný CI systém, který testuje pull
requesty, jen se bude lišit spouštěč události, jenž bude push do testovací
větve \cite{gh_push_trigger}. Po úspěšném sestavení balíčku se testovací 
větev spojí s hlavní větví a úplně se tak vyhneme nutnosti vytvářet pull
request. V případě neúspěchu bude kontaktován správce uvedený v metadatech
balíčku.


\subsection{Issue repozitář}

Pro sledování neaktuálních balíčků existuje separátní issue repozitář, kde
jsou vedeny záznamy o jednotlivých neúspěšně aktualizovaných a tedy neaktuálních
balíčcích. Tento repozitář umožní hlavnímu správci větší přehled o aktuálnosti
balíčků v repozitáři. Záznamy jsou vytvářeny automatizovaně při neúspěšné automatizované
aktualizaci balíčku a stejně tak jsou automatizovaně mazány v případě, že správce
balíček sám aktualizuje.


\subsection{Úložiště a build systém}

Systém separátně využije build systému, který sestavuje balíčky na základě
poskytnutých zdrojů a následně takto sestavený balíček vrátí, aby mohl být
uložen do nějaké databáze, odkud bude dostupný uživatelům k instalaci.
Implementace vlastního build systému je komplikovaný úkol a proto budou
databáze i build systém odděleny od GutHub repozitáře kvůli možnosti
využití již existujícího řešení. Komunikace mezi těmito systémy bude
zajištěna skrze GitHub Actions. Výběr build systému a k tomu vhodného
úložiště jsem diskutoval v následující kapitole.


% ==================================================================================== %
% ==================================================================================== %


\chapter{Metody}

Možnosti, jak zabalit software a následně jej jakkoliv spravovat je
jistě spousta. Pro náš účel jsou vhodné metody, které jsou uživatelsky
přívětivé, software jde s jejich pomocí snadno zabalit a v neposlední řadě
je důležité, aby nad zabaleným softwarem měl plnou kontrolu jeho správce,
který může libovolně a rychle přidávat nebo odstraňovat software. Kvůli poslední
zmíněné vlastnosti tak odpadávají možnosti jako snažit se zabalit software
do nějaké distribuce, poněvadž zabalit software do oficiálních repozitářů
distribucí je typicky časově velmi nákladná činnost.

% ------------------------------------------------------------------------------------ %

\section{Copr}

% ------------------------------------------------------------------------------------ %

\section{Flatpak}


\subsection{Repozitář}

Samotný Flatpak není takovou službou, jak je tomu u Copru, repozitář by bylo
nutno zprovoznit za pomocí zdrojů, které by poskytnula sama fakulta.


% ==================================================================================== %
% ==================================================================================== %


\chapter{Srovnání metod}

V této kapitole jsem se zaměřil na srovnání použitých metod z předchozí
kapitoly. Srovnával jsem na základě předem definovaných kritérií, zvážil
jejich výhody i nevýhody a následně vybral vhodnější metodu k implementaci
řešení.

% ------------------------------------------------------------------------------------ %

\section{Kritéria}

Pro výběr lepší metody je nutné zvolit kritéria, která budou zřetelně hodnotit
různé vlastnosti srovnávané metody. Zvolil jsem čtyři důležitá kritéria:

\paragraph{Nároky na údržbu}
Vzhledem k faktu, že správcem celého systému má být vyučující, musí být správa
celého systému co nejméně časově náročná, ideálně bezúdržbová. Důraz bude
kladen hlavně na části systému, které potenciálně vyžadují pravidelnou
údržbu a také na části, které mohou selhat a bude třeba je znovu zprovoznit
silami fakulty.

\paragraph{Nároky na zdroje}
Zvolené řešení může potřebovat ke svému hladkému provozu nějaké zdroje, které
mohou být například: úložiště, výpočetní výkon nebo potřeba dalšího správce,
který by řešil problémy typu aktualizace celého systému a byl by schopen zprovoznit
systém v případě, že by cokoliv selhalo. Zdroje fakulty samozřejmě nejsou neomezené
a tedy řešení by mělo využívat dostupné zdroje fakulty střídmě a ideálně bez potřeby
správce navíc.

\paragraph{Uživatelská přívětivost}
V ideálním případě by si uživatel měl o nainstalování softwaru na své zařízení
říct jedním příkazem nebo pár kliknutími. Pokud by bylo komplikované software
nainstalovat na zařízení uživatele, ztrácelo by to smysl.

\paragraph{Proveditelnost}
A v neposlední řadě bych měl být schopen celý systém se zvolenou metodou 
implementovat v rozumném časovém rozmezí, které je hodno bakalářské práci 
a to do takové podoby, ve které bude systém použitelný a nebo by mohla být
nad systémem připravena další bakalářská práce, která by celý systém dokončila
do použitelného stavu. V případě extrémní složitosti a časové náročnosti při 
návrhu a implementaci systému jej nemá smysl dělat, jelikož by v tomto případě
nemusel také být nikdy dokončen.

 \paragraph{Spolehlivost}

% ------------------------------------------------------------------------------------ %

\section{Copr}

popis copru

% ------------------------------------------------------------------------------------ %

\section{Flatpak}

popis flatpaku

% ------------------------------------------------------------------------------------ %

\section{Výběr vhodné metody}

sekce o srovnani flatpaku s coprem a vyber viteze


% ==================================================================================== %
% ==================================================================================== %


\chapter{Implementace}

TODO tady budou detaily o implementaci a bla bla


% ==================================================================================== %
% ==================================================================================== %


\chapter*{Závěr}
\markright{\textsc{Závěr}}
\addcontentsline{toc}{chapter}{Závěr}

TODO na konci bakalarky

\printbibliography[heading=bibintoc] %% Print the bibliography.

\end{document}
